\documentclass{dlproblemset}
\usepackage{tikz}
\tikzset{every picture/.style={line width=0.75pt}}

\title{Redes Convolucionais}
\subtitle{Lista de Exercícios}

\begin{document}

\begin{enumerate}[I)]

    \item A rede convolucional LeNet-5 ilustrada na Figura \ref{fig:lenet5} foi desenvolvida por Yann LeCun e publicada em 1998. Sobre a LeNet-5, responda as perguntas abaixo.

    \begin{figure}[H]
        \centering
        \includegraphics[width=0.9\linewidth]{problems/02-cnn/images/LeNET5.jpeg}
        \caption{LeNet-5}
        \label{fig:lenet5}
    \end{figure}
\begin{enumerate}[a)]
    \item Qual o motivo do número 5 no nome dessa rede?
    \item Quais os tamanhos das convoluções usadas?
    \item Qual o tamanho do campo receptivo de uma unidade dos mapas de ativação 14x14 gerados após a camada de \textit{average pooling} (\textit{subsampling})?
    %\item Qual o tamanho do campo receptivo de uma unidade dos mapas de ativação 5x5 gerados após a camada de \textit{average pooling} (\textit{subsampling})?
    \item Quantos parâmetros treináveis tem essa rede?
    \item A maior parte dos parâmetros treináveis está localizada nas camadas convolucionais ou nas totalmente conectadas?
\end{enumerate}

    \item O que é o \textit{stride} em uma camada convolucional?
    \item O que é o \textit{dilation} em uma camada convolucional?
    \item O que é uma convolução transposta?
    \item Por que dizemos que camadas de agrupamento do tipo \textit{Max Pooling} e \textit{Average Pooling} fazem \textit{subsampling} dos mapas de ativação?
    \item Por que usualmente omitimos o número de canais em uma convolução 2D?
    \item Considerando uma convolução 2D 3x3 que conecta dois mapas de ativação com tamanho 224x224 com $C_1$ e $C_2$ canais respectivamente. Como você poderia inicializar os pesos dessa camada usando a inicialização He?

    \item Dada a matriz de input abaixo, e uma rede com 2 camadas convolucionais, calcule a saída da rede com as configurações a seguir. Para a primeira camada convolucional: Padding = 0, Stride = 1; Kernel = $\mathbf{\Theta}_1$; Ativação linear na saída. Para a segunda camada convolucional: Padding = 1; Stride = 1; Kernel = $\mathbf{\Theta}_2$; Ativação ReLU na saída. Considere \textit{dilation} = 1, ou seja, o \textit{kernel} denso.

    $$\mathbf{X} = \begin{bmatrix}
    7 & 8 & 9 \\
    6 & 5 & 4 \\
    1 & 2 & 3
    \end{bmatrix}
    \qquad
    \mathbf{\Theta}_1 = \begin{bmatrix}
    1 & 0 & -1 \\
    1 & 0 & -1 \\
    1 & 0 & -1
    \end{bmatrix}
    \qquad
    \mathbf{\Theta}_2 = \begin{bmatrix}
    1 & 1 & 1 \\
    0 & 0 & 0 \\
    -1 & -1 & -1
    \end{bmatrix}$$


    \item Dadas as alternativas abaixo, qual das alternativas representa o Kernel mais indicado para identificação de bordas horizontais em uma imagem?
\begin{enumerate}[a)]
    \item $$\begin{bmatrix}
    -1 & -1 & -1 \\
    -1 & 8 & -1 \\
    -1 & -1 & -1
    \end{bmatrix}$$
    \item $$\begin{bmatrix}
    -1 & 0 & 1 \\
    -2 & 0 & 2 \\
    -1 & 0 & 1
    \end{bmatrix}$$
    \item $$\begin{bmatrix}
    0 & -1 & 0 \\
    -2 & 6 & -2 \\
    0 & -1 & 0
    \end{bmatrix}$$
    \item $$\begin{bmatrix}
    -1 & -2 & -1 \\
    0 & 0 & 0 \\
    1 & 2 & 1
    \end{bmatrix}$$
\end{enumerate}


    \item Dada uma imagem de entrada de $32\times32$ pixels, e duas camadas convolucionais, qual o campo receptivo na imagem original que é representado por um único pixel na saída final? Considere: Kernel para as duas camadas ($\mathbf{\Theta}_1 = \mathbf{\Theta}_2$) tem tamanho $5\times5$; Stride = 1; \textit{Dilation} = 2.

    \item Uma imagem de entrada $64 \times 64$ passa, em sequência, pelas seguintes camadas (todas quadradas):

    \begin{enumerate}[(i)]
        \item convolução $K = 5$, \textit{stride} $S = 1$, \textit{padding} $P = 0$;
        \item \textit{max pooling} $2 \times 2$, \textit{stride} $S = 2$, \textit{padding} $P = 0$;
        \item convolução $K = 3$, \textit{stride} $S = 1$, \textit{padding} $P = 1$;
        \item \textit{max pooling} $2 \times 2$, \textit{stride} $S = 2$, \textit{padding} $P = 0$;
        \item convolução $K = 3$, \textit{stride} $S = 2$, \textit{padding} $P = 1$.
    \end{enumerate}

    Qual a dimensão espacial (em \textit{pixels}) do mapa de ativação (\textit{feature map}) na saída da camada (v)?

    \begin{enumerate}[a)]
        \item $9 \times 9$
        \item $7 \times 7$
        \item $16 \times 16$
        \item $15 \times 15$
        \item $6 \times 6$
        \item $8 \times 8$
        \item $4 \times 4$
        \item $30 \times 30$
    \end{enumerate}

    \item Avalie as cinco afirmações a seguir sobre redes convolucionais:

    \textbf{I.} O número de parâmetros de uma camada convolucional independe das dimensões espaciais $H \times W$ da entrada, dependendo apenas de $K$, $C_{\text{in}}$, $C_{\text{out}}$ (e dos vieses).

    \textbf{II.} A convolução é equivariante à translação, enquanto o \textit{pooling} (por exemplo, \textit{max pooling}) introduz um grau de invariância a pequenas translações locais.

    \textbf{III.} Durante a inferência, a \textit{batch normalization} normaliza cada ativação usando a média e a variância calculadas sobre o próprio \textit{mini-batch} de teste corrente.

    \textbf{IV.} Duas convoluções $3 \times 3$ empilhadas têm o mesmo campo receptivo de uma única convolução $5 \times 5$, porém usam \emph{mais} parâmetros do que a $5 \times 5$.

    \textbf{V.} Uma convolução $1 \times 1$ preserva as dimensões espaciais e serve, entre outras coisas, para alterar o número de canais (combinação linear entre canais em cada posição).

    \begin{enumerate}[a)]
        \item Apenas I e V estão corretas.
        \item Apenas I, II, IV e V estão corretas.
        \item Apenas II, III e V estão corretas.
        \item Apenas I, II e V estão corretas.
        \item Apenas I, II e III estão corretas.
        \item Apenas I, III e IV estão corretas.
        \item Todas as afirmações estão corretas.
        \item Apenas II e V estão corretas.
    \end{enumerate}

    \item Considere três camadas convolucionais em sequência, com os valores de \textit{kernel} $K$, \textit{stride} $S$ e dilatação $D$ indicados abaixo:

    \begin{itemize}
        \item camada 1: $K_1 = 3$, $S_1 = 1$, $D_1 = 2$;
        \item camada 2: $K_2 = 3$, $S_2 = 2$, $D_2 = 1$;
        \item camada 3: $K_3 = 3$, $S_3 = 1$, $D_3 = 1$.
    \end{itemize}

    Qual a dimensão (em \textit{pixels}) do campo receptivo de uma unidade na saída da camada 3?

    \begin{enumerate}[a)]
        \item $9 \times 9$
        \item $11 \times 11$
        \item $13 \times 13$
        \item $7 \times 7$
        \item $15 \times 15$
        \item $5 \times 5$
        \item $17 \times 17$
        \item $8 \times 8$
    \end{enumerate}

    \item Considere uma CNN com três camadas convolucionais em sequência, todas \textbf{sem} \textit{padding}. A camada 1 usa \textit{kernel} $3\times 3$ com \textit{stride} $S_1 = 1$; a camada 2 usa \textit{kernel} $3\times 3$ com \textit{stride} $S_2 = 2$; a camada 3 usa \textit{kernel} $3\times 3$ com \textit{stride} $S_3 = 1$. Qual o campo receptivo (lado, em \textit{pixels}) de uma unidade na saída da camada 3?
    \begin{enumerate}[a)]
        \item $5$
        \item $6$
        \item $7$
        \item $9$
        \item $11$
        \item $13$
        \item $15$
        \item $17$
    \end{enumerate}

    \item Um colega propõe usar um MLP para classificar imagens RGB de dimensão $64 \times 64 \times 3$, achatando a imagem em um vetor de entrada e usando uma primeira camada oculta totalmente conectada com $256$ unidades.

    \begin{enumerate}[a)]
        \item Compute o número de parâmetros treináveis (pesos e vieses) dessa primeira camada totalmente conectada e compare-o com o número de parâmetros de uma camada convolucional $3 \times 3$ com $32$ filtros (com viés), \textit{stride} $1$ e \textit{padding} $1$, aplicada à mesma imagem. Mostre o cálculo dos dois valores.
        \item Além da mudança no número de parâmetros, aponte \textbf{duas} propriedades estruturais da convolução que a tornam adequada para imagens, explicando em uma ou duas frases por que cada uma delas ajuda nessa tarefa.
    \end{enumerate}

\end{enumerate}

\newpage
\subsection*{Gabarito}
\color{blue}

\begin{enumerate}[I)]
\item
\begin{enumerate}[a.]
    \item O número 5 indica o número de camadas com parâmetros treináveis. São 2 convoluções e 3 totalmente conectadas.

    \item Na primeira camada convolucional foram convoluções 5x5 com 6 canais de saída; Na segunda foram convoluções 5x5 com 16 canais de saída.

    \item Tamanho 6x6. Uma convolução 5x5 tem campo receptivo de tamanho 5x5, se depois dele fizermos um \textit{average pooling} 2x2 teremos \textit{receptive field} de 6x6.

    %\item Tamanho 16x16. Para responder a essa questão você deve perceber que ao aplicar uma conv após fazer pooling 2x2 com stride 2 o campo perceptivo dobra de tamanho.

    \item Contando as camadas com parâmetros treináveis e assumindo que C3 está totalmente conectada aos 6 mapas de S2: primeira camada convolucional $5\times5\times1\times6 + 6 = 156$; segunda camada convolucional $5\times5\times6\times16 + 16 = 2416$; terceira camada $400\times120 + 120 = 48\,120$; quarta camada $120\times84 + 84 = 10\,164$; quinta camada $84\times10 + 10 = 850$. Total $61\,706$.

    Vale registrar que a LeNet-5 do artigo de 1998 tem cerca de $60\,840$ parâmetros, número que o próprio artigo arredonda para 60 mil. A diferença em relação aos $61\,706$ acima vem de três escolhas de modelagem que a leitura moderna da arquitetura abandonou:
    \begin{itemize}
        \item \textbf{C3 com conectividade parcial.} No artigo, cada um dos 16 mapas de C3 recebe apenas um subconjunto dos 6 mapas de S2 (Tabela I do artigo), o que dá 60 \textit{kernels} $5\times5$ em vez de 96, ou seja $60\times25 + 16 = 1516$ parâmetros em vez de $2416$. A motivação era dupla: reduzir o custo computacional e quebrar a simetria entre os mapas, forçando cada um a extrair características diferentes.
        \item \textbf{Camada de saída RBF.} No artigo a saída não é uma camada linear com \textit{bias}, e sim uma camada RBF que mede a distância euclidiana a 10 vetores-protótipo de 84 dimensões, ou seja $84\times10 = 840$ pesos em vez de $850$. Esses protótipos foram definidos à mão e mantidos fixos, então contá-los como treináveis é uma convenção.
        \item \textbf{Subamostragem com parâmetros.} S2 e S4 não são \textit{average pooling} puro: no artigo cada mapa tem um coeficiente multiplicativo e um \textit{bias} treináveis, somando $6\times2 + 16\times2 = 44$ parâmetros que a contagem acima não inclui, porque hoje tratamos \textit{pooling} como operação sem parâmetros.
    \end{itemize}
    Fechando a conta: $61\,706 - 900 - 10 + 44 = 60\,840$.

    \item Nas camadas totalmente conectadas.
\end{enumerate}
\item É o quanto o filtro translada entre cada uma das aplicações. Por padrão em convoluções usamos 1, o que significa que o kernel é transladado de uma posição após cada aplicação do filtro.

\item É um parâmetro usado para aumentar o campo receptivo da convolução. Na prática tornamos o kernel esparso completando com zeros.

\item É uma convolução usada para fazer \textit{upsampling} na imagem. É uma forma de aprender a função de \textit{upsampling}.

\item Pois normalmente usamos essas camadas para diminuir a resolução espacial dos mapas de ativação.

\item Pois usualmente o número de canais da convolução é igual ao número de canais dos mapas de ativação.

\item A inicialização He leva em conta o número de parâmetros que contribuem para cada ativação. Como nesse caso estamos falando de $C_1x3x3$ pesos, podemos inicializar com uma normal com média zero e variância $\frac{2}{9C_1}$

\item Resolução Input 3x3 % resolução questão 8 VIII
    Matrizes Iniciais
    $$
    \mathbf{X} = \begin{bmatrix}
    7 & 8 & 9 \\
    6 & 5 & 4 \\
    1 & 2 & 3
    \end{bmatrix}
    \qquad
    \mathbf{\Theta}_1 = \begin{bmatrix}
    1 & 0 & -1 \\
    1 & 0 & -1 \\
    1 & 0 & -1
    \end{bmatrix}
    \qquad
    \mathbf{\Theta}_2 = \begin{bmatrix}
    1 & 1 & 1 \\
    0 & 0 & 0 \\
    -1 & -1 & -1
    \end{bmatrix}
    $$

    Camada 1 (Conv1)
    Input $3\times3$, Kernel $3\times3$, $p=0$, $s=1 \implies$ Saída $1\times1$.
    $$
    \mathbf{C}_1 = \sum(\mathbf{X} \cdot \mathbf{\Theta}_1) + B_1 \quad (B_1 = 0)
    $$
    $$
    \mathbf{C}_1 = (7\cdot1 + 8\cdot0 + 9\cdot-1) + (6\cdot1 + 5\cdot0 + 4\cdot-1) + (1\cdot1 + 2\cdot0 + 3\cdot-1)
    $$
    $$
    \mathbf{C}_1 = (-2) + (2) + (-2) = -2
    $$
    $$
    \mathbf{C}_1 = \begin{bmatrix} -2 \end{bmatrix}
    $$

    Camada 2 (Conv2)
    Input $1\times1$, Kernel $3\times3$, $p=1$, $s=1 \implies$ Saída $1\times1$.
    Primeiro, aplicamos padding na saída de $\mathbf{C}_1$:
    $$
    \mathbf{C}_{1, padded} = \begin{bmatrix}
    0 & 0 & 0 \\
    0 & -2 & 0 \\
    0 & 0 & 0
    \end{bmatrix}
    $$
    Em seguida, aplicamos $\mathbf{\Theta}_2$ e a ReLU. O enunciado não define \textit{bias} em nenhuma das duas camadas, então $B_1 = B_2 = 0$. O elemento central de $\mathbf{\Theta}_2$ é $0$ e a única posição não nula de $\mathbf{C}_{1,\,padded}$ é justamente a central, logo a soma se anula:
    $$
    \text{Soma Conv} = \sum(\mathbf{C}_{1,\,padded} \cdot \mathbf{\Theta}_2) = (-2) \cdot 0 = 0
    $$
    $$
    \text{Pré-ativação} = \text{Soma Conv} + B_2 = 0 + 0 = 0
    $$
    $$
    \text{Output} = \text{ReLU}(0) = 0
    $$

    Resultado Final
    $$
    \text{Output} = \begin{bmatrix} 0 \end{bmatrix}
    $$

\item Alternativa (d). É o operador de Sobel na direção vertical, $\partial/\partial y$: as linhas superior e inferior têm sinais opostos, então a resposta é máxima onde a intensidade muda de cima para baixo, que é o que caracteriza uma borda horizontal. As demais: (a) é o Laplaciano, um filtro passa-alta isotrópico que responde a bordas em qualquer orientação, sem selecionar direção; (b) é o Sobel na direção horizontal, $\partial/\partial x$, com as colunas de sinais opostos, então detecta bordas \textbf{verticais}; (c) é um filtro de realce, também de soma nula, com pesos maiores na direção horizontal, mas que responde às duas orientações e serve para acentuar detalhes, não para isolar bordas horizontais. % resolução 9 IX

   \item A primeira camada de convolução com kernel $5\times5$ e dilatação $2$, resulta em um \textit{feature map} com campo receptivo de $9\times9$ sobre a imagem de entrada. A segunda camada aplica a mesma configuração de kernel sobre o \textit{feature map} da primeira camada de convolução. Portanto, o campo receptivo final é o campo de visão da camada anterior ($9\times9$) mais a expansão adicional do seu próprio kernel. Resultando em um campo receptivo de $17\times17$. % questão X

\item Alternativa (f). Aplicando $H' = \lfloor (H - K + 2P)/S \rfloor + 1$ camada a camada, partindo de $64$:
\begin{align*}
    \text{(i) conv } K{=}5, S{=}1, P{=}0:&\ (64 - 5)/1 + 1 = 60 \\
    \text{(ii) pool } 2, S{=}2:&\ (60 - 2)/2 + 1 = 30 \\
    \text{(iii) conv } K{=}3, S{=}1, P{=}1:&\ (30 - 3 + 2)/1 + 1 = 30 \\
    \text{(iv) pool } 2, S{=}2:&\ (30 - 2)/2 + 1 = 15 \\
    \text{(v) conv } K{=}3, S{=}2, P{=}1:&\ \lfloor (15 - 3 + 2)/2 \rfloor + 1 = \lfloor 14/2 \rfloor + 1 = 8
\end{align*}
Dimensão de saída $8 \times 8$. A alternativa (b) $7 \times 7$ esquece o ``$+1$'' (ou erra o piso) na camada (v); (d) $15 \times 15$ para na camada (iv); (c) $16 \times 16$ trata a camada (v) como preservando a dimensão; (a), (e), (g) e (h) resultam de outros erros de aplicação da fórmula. % questão XI

\item Alternativa (d). I (V): a convolução compartilha pesos, então a contagem depende de $K$, $C_{\text{in}}$, $C_{\text{out}}$ (e vieses), e independe de $H\times W$. II (V): a convolução é equivariante à translação e o \textit{pooling} adiciona invariância local. III (F): na inferência a \textit{batch norm} usa as estatísticas móveis ($\mu_{\text{mov}}, \sigma_{\text{mov}}$) acumuladas no treino, e não as do \textit{mini-batch} de teste (caso contrário a saída de uma amostra dependeria das demais do \textit{batch}). IV (F): duas $3\times 3$ têm de fato o mesmo campo receptivo ($5$) de uma $5\times 5$, porém com \emph{menos} parâmetros ($2\cdot 9 = 18 < 25$) e mais não-linearidades; ``mais parâmetros'' é falso. V (V): a convolução $1\times 1$ preserva o espaço e recombina canais. Verdadeiras: I, II e V. As alternativas (b) e (f) incluem IV (falsa); (c) e (e) incluem III (falsa); (a) omite II; (h) omite I; (g) marca tudo. % questão XII

\item Alternativa (b). \textit{Kernel} efetivo da camada 1: $K_{\text{eff}} = D_1(K_1 - 1) + 1 = 2\cdot 2 + 1 = 5$ (camadas 2 e 3 têm $D = 1$, logo $K_{\text{eff}} = K$). Aplicando a recursão $RF_l = RF_{l-1} + (K_{\text{eff},l} - 1)\prod_{i<l} S_i$ com $RF_0 = 1$:
\begin{align*}
    RF_1 &= 1 + (5 - 1)\cdot 1 = 5 \\
    RF_2 &= 5 + (3 - 1)\cdot S_1 = 5 + 2\cdot 1 = 7 \\
    RF_3 &= 7 + (3 - 1)\cdot (S_1 S_2) = 7 + 2\cdot(1\cdot 2) = 11
\end{align*}
Campo receptivo $11 \times 11$. A alternativa (a) $9 \times 9$ ignora a dilatação na camada 1 (usa $K_{\text{eff}} = 3$: $3 \to 5 \to 9$); (d) $7 \times 7$ para no $RF_2$; (f) $5 \times 5$ para no $RF_1$; (g) $17 \times 17$ é o resultado de \emph{duas} camadas $5\times 5$ dilatadas ($D=2$) sem \textit{stride}; (c), (e) e (h) usam o \textit{stride} da própria camada no produto ou cometem outros erros. % questão XIII

\item Alternativa (d). Aplicando a recursão $RF_l = RF_{l-1} + (K_l - 1)\prod_{i<l} S_i$ com $RF_0 = 1$ e todos os \textit{kernels} $3\times 3$:
\begin{align*}
    RF_1 &= 1 + (3-1)\cdot 1 = 3 \\
    RF_2 &= 3 + (3-1)\cdot S_1 = 3 + 2\cdot 1 = 5 \\
    RF_3 &= 5 + (3-1)\cdot (S_1 S_2) = 5 + 2\cdot(1\cdot 2) = 5 + 4 = 9
\end{align*}
O campo receptivo é $9 \times 9$, logo o lado é $9$. A alternativa (c) $7$ trata todos os \textit{strides} como $1$ ($1 + 2 + 2 + 2 = 7$), ignorando o produto acumulado; (a) $5$ para no $RF_2$, esquecendo a terceira camada; (e) $11$, (f) $13$, (g) $15$ e (h) $17$ usam o \textit{stride} da própria camada no produto (incluindo $i = l$) ou acumulam o \textit{stride} $2$ em camadas erradas; (b) $6$ é um número sem correspondência com a recursão. % questão XIV

\item MLP \textit{versus} convolução para imagens. % questão XV
\begin{enumerate}[a.]
    \item Camada totalmente conectada: a entrada achatada tem $64 \cdot 64 \cdot 3 = 12288$ componentes; com $256$ unidades,
    \[
    12288 \times 256 + 256 = 3\,145\,728 + 256 = 3\,145\,984 \text{ parâmetros}.
    \]
    Camada convolucional $3 \times 3$ com $3$ canais de entrada e $32$ filtros:
    \[
    3 \times 3 \times 3 \times 32 + 32 = 864 + 32 = 896 \text{ parâmetros},
    \]
    cerca de $3500$ vezes menos. Note que o total da convolução independe das dimensões espaciais $64 \times 64$ da imagem.

    \item Duas propriedades (quaisquer duas, bem explicadas, respondem ao item):
    \begin{itemize}
        \item \textbf{Compartilhamento de pesos:} o mesmo \textit{kernel} é aplicado em todas as posições da imagem, tornando a operação equivariante à translação: um padrão aprendido em um canto é reconhecido em qualquer outra posição sem parâmetros adicionais.
        \item \textbf{Conectividade local:} cada unidade olha apenas para uma vizinhança pequena, explorando a coesão espacial das imagens (\textit{pixels} próximos são fortemente correlacionados), e o empilhamento de camadas amplia o campo receptivo de forma hierárquica.
        \item (Também aceito) \textbf{Independência do tamanho da entrada:} o número de parâmetros depende só de \textit{kernel} e canais, permitindo aplicar a mesma camada a imagens de resoluções diferentes.
    \end{itemize}
\end{enumerate}

\end{enumerate}
\end{document}
